\documentclass[11pt]{article}
\usepackage[right=1in,left=1in,top=0.85in,bottom=0.85in]{geometry}
%% ============================================================
%% Common notation for all econirl primer documents.
%%
%% Usage: place %% ============================================================
%% Common notation for all econirl primer documents.
%%
%% Usage: place %% ============================================================
%% Common notation for all econirl primer documents.
%%
%% Usage: place \input{../common/notation} in your preamble
%%        (before \begin{document}), then call \commonnotation
%%        inside the Model and Notation section.
%%
%% This file provides:
%%   1. Shared LaTeX macros (operators, calligraphic sets)
%%   2. \commonnotation command — a DDC notation table
%%   3. Standard package imports used by every primer
%% ============================================================

%% ---------- packages every primer needs ----------
%% Algorithm packages are NOT loaded here because primers differ:
%%   algorithm2e  → SEES, NNES, MCE-IRL, MaxEnt, Deep MCE-IRL, CCP
%%   algorithm    → AIRL, AAIRL
%% Each primer loads its own algorithm package after this \input.
\usepackage{amsmath,amssymb}
\usepackage{booktabs}
\usepackage[round]{natbib}
\usepackage{hyperref}
\hypersetup{colorlinks,citecolor=blue,filecolor=blue,linkcolor=blue,urlcolor=blue}
\usepackage{float}
\usepackage{microtype}
\usepackage{placeins}

%% ---------- shared macros ----------
\usepackage{amsthm}
\newcommand{\E}{\mathbb{E}}
\newcommand{\R}{\mathbb{R}}
\newcommand{\calS}{\mathcal{S}}
\newcommand{\calA}{\mathcal{A}}
\newcommand{\calH}{\mathcal{H}}
\newcommand{\calD}{\mathcal{D}}
\DeclareMathOperator*{\argmin}{arg\,min}
\DeclareMathOperator*{\argmax}{arg\,max}

%% ---------- theorem environments ----------
\newtheorem{proposition}{Proposition}
\newtheorem{remark}{Remark}
\newtheorem{definition}{Definition}

%% ---------- code listings ----------
\usepackage{listings}
\lstset{
  language=Python,
  basicstyle=\small\ttfamily,
  keywordstyle=\bfseries,
  breaklines=true,
  frame=single,
  xleftmargin=1em,
  numbers=none,
}

%% ---------- common DDC notation table ----------
\newcommand{\commonnotation}{%
\begin{table}[H]
\centering\small
\caption{Common DDC notation shared across all primers.}
\label{tab:common_notation}
\begin{tabular*}{\textwidth}{@{\extracolsep{\fill}} l l l}
\toprule
Symbol & Name & Definition \\
\midrule
$s \in \calS$ & State & Discrete or continuous \\
$a \in \calA$ & Action & Finite: $|\calA| = J$ \\
$s'$ & Next state & Drawn from $P(s' \mid s, a)$ \\
$\beta \in [0,1)$ & Discount factor & \\
$\sigma > 0$ & Scale parameter & Logit noise scale ($= 1$ by convention) \\
$\theta \in \Theta \subset \R^{d_\theta}$ & Structural parameters & Unknown \\
$\phi(s,a)$ & Feature vector & Known functions of $(s,a)$, dimension $K$ \\
$u(s,a;\theta)$ & Flow utility & $\theta^\top \phi(s,a)$ (linear-in-parameters) \\
$P(s' \mid s, a)$ & Transition kernel & Known or estimated from data \\
$Q(s,a)$ & Choice-specific value & $u(s,a;\theta) + \beta \sum_{s'} P(s'|s,a)\, V(s')$ \\
$V(s)$ & Integrated value & $\sigma \log \sum_a \exp(Q(s,a)/\sigma)$ \\
$\pi(a \mid s)$ & Choice probability & $\exp\bigl((Q(s,a) - V(s))/\sigma\bigr)$ \\
$\tau$ & Trajectory & $(s_0, a_0, s_1, a_1, \ldots)$ from one agent \\
$\calD$ & Panel data & Collection of $N$ observed trajectories \\
\bottomrule
\end{tabular*}
\end{table}%
}
 in your preamble
%%        (before \begin{document}), then call \commonnotation
%%        inside the Model and Notation section.
%%
%% This file provides:
%%   1. Shared LaTeX macros (operators, calligraphic sets)
%%   2. \commonnotation command — a DDC notation table
%%   3. Standard package imports used by every primer
%% ============================================================

%% ---------- packages every primer needs ----------
%% Algorithm packages are NOT loaded here because primers differ:
%%   algorithm2e  → SEES, NNES, MCE-IRL, MaxEnt, Deep MCE-IRL, CCP
%%   algorithm    → AIRL, AAIRL
%% Each primer loads its own algorithm package after this \input.
\usepackage{amsmath,amssymb}
\usepackage{booktabs}
\usepackage[round]{natbib}
\usepackage{hyperref}
\hypersetup{colorlinks,citecolor=blue,filecolor=blue,linkcolor=blue,urlcolor=blue}
\usepackage{float}
\usepackage{microtype}
\usepackage{placeins}

%% ---------- shared macros ----------
\usepackage{amsthm}
\newcommand{\E}{\mathbb{E}}
\newcommand{\R}{\mathbb{R}}
\newcommand{\calS}{\mathcal{S}}
\newcommand{\calA}{\mathcal{A}}
\newcommand{\calH}{\mathcal{H}}
\newcommand{\calD}{\mathcal{D}}
\DeclareMathOperator*{\argmin}{arg\,min}
\DeclareMathOperator*{\argmax}{arg\,max}

%% ---------- theorem environments ----------
\newtheorem{proposition}{Proposition}
\newtheorem{remark}{Remark}
\newtheorem{definition}{Definition}

%% ---------- code listings ----------
\usepackage{listings}
\lstset{
  language=Python,
  basicstyle=\small\ttfamily,
  keywordstyle=\bfseries,
  breaklines=true,
  frame=single,
  xleftmargin=1em,
  numbers=none,
}

%% ---------- common DDC notation table ----------
\newcommand{\commonnotation}{%
\begin{table}[H]
\centering\small
\caption{Common DDC notation shared across all primers.}
\label{tab:common_notation}
\begin{tabular*}{\textwidth}{@{\extracolsep{\fill}} l l l}
\toprule
Symbol & Name & Definition \\
\midrule
$s \in \calS$ & State & Discrete or continuous \\
$a \in \calA$ & Action & Finite: $|\calA| = J$ \\
$s'$ & Next state & Drawn from $P(s' \mid s, a)$ \\
$\beta \in [0,1)$ & Discount factor & \\
$\sigma > 0$ & Scale parameter & Logit noise scale ($= 1$ by convention) \\
$\theta \in \Theta \subset \R^{d_\theta}$ & Structural parameters & Unknown \\
$\phi(s,a)$ & Feature vector & Known functions of $(s,a)$, dimension $K$ \\
$u(s,a;\theta)$ & Flow utility & $\theta^\top \phi(s,a)$ (linear-in-parameters) \\
$P(s' \mid s, a)$ & Transition kernel & Known or estimated from data \\
$Q(s,a)$ & Choice-specific value & $u(s,a;\theta) + \beta \sum_{s'} P(s'|s,a)\, V(s')$ \\
$V(s)$ & Integrated value & $\sigma \log \sum_a \exp(Q(s,a)/\sigma)$ \\
$\pi(a \mid s)$ & Choice probability & $\exp\bigl((Q(s,a) - V(s))/\sigma\bigr)$ \\
$\tau$ & Trajectory & $(s_0, a_0, s_1, a_1, \ldots)$ from one agent \\
$\calD$ & Panel data & Collection of $N$ observed trajectories \\
\bottomrule
\end{tabular*}
\end{table}%
}
 in your preamble
%%        (before \begin{document}), then call \commonnotation
%%        inside the Model and Notation section.
%%
%% This file provides:
%%   1. Shared LaTeX macros (operators, calligraphic sets)
%%   2. \commonnotation command — a DDC notation table
%%   3. Standard package imports used by every primer
%% ============================================================

%% ---------- packages every primer needs ----------
%% Algorithm packages are NOT loaded here because primers differ:
%%   algorithm2e  → SEES, NNES, MCE-IRL, MaxEnt, Deep MCE-IRL, CCP
%%   algorithm    → AIRL, AAIRL
%% Each primer loads its own algorithm package after this \input.
\usepackage{amsmath,amssymb}
\usepackage{booktabs}
\usepackage[round]{natbib}
\usepackage{hyperref}
\hypersetup{colorlinks,citecolor=blue,filecolor=blue,linkcolor=blue,urlcolor=blue}
\usepackage{float}
\usepackage{microtype}
\usepackage{placeins}

%% ---------- shared macros ----------
\usepackage{amsthm}
\newcommand{\E}{\mathbb{E}}
\newcommand{\R}{\mathbb{R}}
\newcommand{\calS}{\mathcal{S}}
\newcommand{\calA}{\mathcal{A}}
\newcommand{\calH}{\mathcal{H}}
\newcommand{\calD}{\mathcal{D}}
\DeclareMathOperator*{\argmin}{arg\,min}
\DeclareMathOperator*{\argmax}{arg\,max}

%% ---------- theorem environments ----------
\newtheorem{proposition}{Proposition}
\newtheorem{remark}{Remark}
\newtheorem{definition}{Definition}

%% ---------- code listings ----------
\usepackage{listings}
\lstset{
  language=Python,
  basicstyle=\small\ttfamily,
  keywordstyle=\bfseries,
  breaklines=true,
  frame=single,
  xleftmargin=1em,
  numbers=none,
}

%% ---------- common DDC notation table ----------
\newcommand{\commonnotation}{%
\begin{table}[H]
\centering\small
\caption{Common DDC notation shared across all primers.}
\label{tab:common_notation}
\begin{tabular*}{\textwidth}{@{\extracolsep{\fill}} l l l}
\toprule
Symbol & Name & Definition \\
\midrule
$s \in \calS$ & State & Discrete or continuous \\
$a \in \calA$ & Action & Finite: $|\calA| = J$ \\
$s'$ & Next state & Drawn from $P(s' \mid s, a)$ \\
$\beta \in [0,1)$ & Discount factor & \\
$\sigma > 0$ & Scale parameter & Logit noise scale ($= 1$ by convention) \\
$\theta \in \Theta \subset \R^{d_\theta}$ & Structural parameters & Unknown \\
$\phi(s,a)$ & Feature vector & Known functions of $(s,a)$, dimension $K$ \\
$u(s,a;\theta)$ & Flow utility & $\theta^\top \phi(s,a)$ (linear-in-parameters) \\
$P(s' \mid s, a)$ & Transition kernel & Known or estimated from data \\
$Q(s,a)$ & Choice-specific value & $u(s,a;\theta) + \beta \sum_{s'} P(s'|s,a)\, V(s')$ \\
$V(s)$ & Integrated value & $\sigma \log \sum_a \exp(Q(s,a)/\sigma)$ \\
$\pi(a \mid s)$ & Choice probability & $\exp\bigl((Q(s,a) - V(s))/\sigma\bigr)$ \\
$\tau$ & Trajectory & $(s_0, a_0, s_1, a_1, \ldots)$ from one agent \\
$\calD$ & Panel data & Collection of $N$ observed trajectories \\
\bottomrule
\end{tabular*}
\end{table}%
}

\usepackage{algorithm}
\usepackage{algorithmic}
\usepackage{titling}
\setlength{\droptitle}{-2.5em}
\posttitle{\par\end{center}\vspace{-1.2em}}
\setlength{\parskip}{0.3em}
\bibliographystyle{abbrvnat}

\title{Primer: Generative Adversarial Imitation Learning (GAIL)\\[3pt]
\large Ho and Ermon (2016)}
\author{econirl package documentation}
\date{}

\begin{document}
\maketitle
\vspace{-0.5em}

%% ============================================================
\section{Context and Problem}
%% ============================================================

Traditional IRL solves a two-step problem: first recover the reward from expert demonstrations, then find the optimal policy for that reward.
Both steps are expensive.
The reward recovery requires solving an MDP at each gradient step, and the policy optimization requires solving a second MDP after IRL converges.
\citet{hoErmon2016} observe that if the goal is only to match the expert's behavior (not to understand the reward), the two steps can be collapsed into one.
Their key insight is that IRL is dual to occupancy measure matching.
An occupancy measure $\rho_\pi(s,a)$ counts how often a policy $\pi$ visits state-action pair $(s,a)$ under the discounted visitation distribution.
Two policies produce the same behavior if and only if they have the same occupancy measure \citep{puterman1994}.
So matching the expert's behavior reduces to matching its occupancy measure, which can be done directly via a GAN \citep{goodfellow2014} without ever recovering a reward.

The tradeoff is clear: GAIL is faster and simpler than IRL but does not produce a reward function.
This means GAIL cannot answer counterfactual questions (``what would the agent do under different dynamics?'') because there is no structural primitive to re-optimize.
For structural estimation and counterfactual analysis, AIRL \citep{fu2018} or MCE-IRL \citep{ziebart2010} are needed.
GAIL serves as a strong behavioral cloning baseline that exploits MDP structure.

%% ============================================================
\section{Model and Notation}
%% ============================================================

\commonnotation

GAIL introduces one additional object: the occupancy measure $\rho_\pi(s,a) = (1-\beta) \sum_{t=0}^\infty \beta^t \Pr(s_t = s, a_t = a \mid \pi)$, which is the discounted frequency of visiting $(s,a)$ under policy $\pi$.
The key equivalence (Proposition 3.1 in \citealt{hoErmon2016}): $\pi \to \rho_\pi$ is a bijection between stationary policies and valid occupancy measures.
Matching $\rho_\pi = \rho_E$ is equivalent to matching $\pi = \pi_E$ in behavioral terms.

%% ============================================================
\section{Theory}
%% ============================================================

\paragraph{IRL as occupancy matching (Proposition 3.2).}
\citet{hoErmon2016} show that the composition $\text{RL} \circ \text{IRL}(\pi_E)$ solves
\begin{equation}
\label{eq:irl_dual}
\min_\pi \; \psi^*(\rho_\pi - \rho_E) - \calH(\pi)
\end{equation}
where $\psi^*$ is the convex conjugate of the cost regularizer $\psi$ used in IRL, and $\calH(\pi)$ is the causal entropy.
Different choices of $\psi$ recover different algorithms.
Setting $\psi$ to a constant gives exact occupancy matching but is impractical.
Setting $\psi$ to an indicator on linear functions gives apprenticeship learning \citep{abbeelNg2004}.

\paragraph{GAIL regularizer (Corollary 3.2.1).}
Choosing $\psi = \psi_{\text{GA}}$, the regularizer induced by a GAN discriminator, gives
\begin{equation}
\label{eq:gail_objective}
\min_\pi \max_D \; \E_{\rho_\pi}[\log D(s,a)] + \E_{\rho_E}[\log(1 - D(s,a))] - \lambda \calH(\pi)
\end{equation}
where $D : \calS \times \calA \to (0,1)$ is a binary classifier.
The inner maximization over $D$ computes the Jensen-Shannon divergence between $\rho_\pi$ and $\rho_E$.
The outer minimization over $\pi$ finds the policy whose occupancy is closest to the expert's.

\paragraph{Why no reward.}
The discriminator $D(s,a)$ is not a reward function.
It is a classifier that distinguishes expert transitions from policy transitions.
At the optimum $D^* = 1/2$ everywhere (expert and policy are indistinguishable), and $D$ carries no information about the underlying reward.
Unlike AIRL, there is no structural decomposition that separates reward from value function.

%% ============================================================
\section{Algorithm}
%% ============================================================

\begin{algorithm}[h!]
\caption{GAIL (Ho \& Ermon 2016, Algorithm 1)}
\label{alg:gail}
\begin{algorithmic}[1]
\REQUIRE Expert demonstrations $\calD_E$, entropy weight $\lambda$
\ENSURE Policy $\pi_\omega$ matching expert behavior

\STATE Initialize discriminator $D_\psi$, policy $\pi_\omega$

\FOR{iteration $= 1, 2, \ldots$}

\STATE Sample expert batch: $\{(s,a)\} \sim \calD_E$
\STATE Sample policy batch: $\{(s,a)\}$ from rollouts of $\pi_\omega$

\STATE \textbf{// Discriminator step}
\STATE $\calL_D = -\E_E[\log D_\psi(s,a)] - \E_\pi[\log(1 - D_\psi(s,a))]$
\STATE $\psi \leftarrow \psi - \alpha_D \nabla_\psi \calL_D$ \hfill \textit{// Adam}

\STATE \textbf{// Policy step}
\STATE Reward signal: $\tilde{r}(s,a) = -\log(1 - D_\psi(s,a))$ \hfill \textit{// or softplus variant}
\STATE Update $\pi_\omega$ to maximize $\E_\pi[\sum_t \beta^t (\tilde{r} + \lambda \calH)]$
\STATE \quad \textbf{Tabular:} soft value iteration on $\tilde{r}$
\STATE \quad \textbf{Neural:} PPO/TRPO on $\tilde{r}$ \citep{schulman2017}

\IF{$D \approx 1/2$ on held-out transitions}
  \STATE \textbf{break} \hfill \textit{// policy matches expert}
\ENDIF

\ENDFOR
\RETURN $\pi_\omega$ \quad {\small (no reward function returned)}
\end{algorithmic}
\end{algorithm}
\FloatBarrier

The reward signal $\tilde{r} = -\log(1 - D)$ encourages the policy to generate transitions that fool the discriminator.
As $D \to 1/2$, the reward becomes uninformative and training stabilizes.
The econirl implementation offers two reward transforms: \texttt{softplus} ($\log(1 + e^{\text{logit}})$, smoother gradients) and \texttt{logit} (raw discriminator output).

%% ============================================================
\section{Results from the Paper}
%% ============================================================

\paragraph{MuJoCo (Table 1).}
On Reacher, HalfCheetah, Hopper, Walker, and Ant, GAIL matches or exceeds the expert's cumulative reward using only 200 expert trajectories.
Behavioral cloning requires 10x more data to achieve the same performance.
Feature-expectation matching (FEM, \citealt{abbeelNg2004}) fails on high-dimensional tasks because the linear feature assumption is too restrictive.

\paragraph{Scalability.}
GAIL scales to high-dimensional continuous state-action spaces via neural discriminators and PPO policy updates.
IRL methods like MaxEnt or MCE-IRL require tabular occupancy computation and do not scale beyond a few hundred states.

\paragraph{Data efficiency.}
GAIL is more data-efficient than BC because it uses the MDP's transition structure.
BC treats each $(s,a)$ pair independently, suffering from covariate shift: errors compound over trajectories because the learner visits states not seen in demonstrations.
GAIL avoids this by rolling out the learned policy and comparing the resulting occupancy to the expert's.

%% ============================================================
\section{Simulation}
%% ============================================================

We test GAIL on a $5 \times 5$ gridworld (25 states, 4 actions, stochastic transitions, $\beta = 0.95$, goal reward at bottom-right).
Expert data: 300 trajectories of 50 steps.
We compare GAIL against behavioral cloning (BC) on two metrics: greedy policy match rate (fraction of states with correct greedy action) and Jensen-Shannon divergence between learned and expert occupancy measures.
Run \texttt{python gail\_run.py} to regenerate Table~\ref{tab:results}.

\begin{table}[h!]
\centering\small
\begin{tabular}{lrr}
\toprule
Metric & BC & GAIL \\
\midrule
% Auto-generated by gail_run.py -- do not edit by hand
Policy match rate & $0.000$ & $0.000$ \\
Occupancy JS divergence & $0.0000$ & $0.0000$ \\
Reward recovered & $No$ & $No$ \\
Time (s) & $---$ & $0.0$ \\

\bottomrule
\end{tabular}
\caption{GAIL vs BC on 25-state gridworld. GAIL should achieve lower occupancy JS divergence (its direct objective) while neither method recovers a reward. Auto-generated by \texttt{gail\_run.py}.}
\label{tab:results}
\end{table}
\FloatBarrier

GAIL's objective~\eqref{eq:gail_objective} directly minimizes the JS divergence between learned and expert occupancy, so it should achieve a lower JS than BC on this metric.
BC matches the expert's conditional action distribution $\pi(a|s)$ at observed states but does not account for the state visitation distribution.
This causes covariate shift: errors at one state push the agent into states underrepresented in demonstrations, where further errors compound.
GAIL avoids this by optimizing over rollouts of the learned policy, so the state distribution under the learned policy enters the objective.

Neither method produces a reward function.
GAIL's discriminator $D(s,a)$ at convergence is $1/2$ everywhere and carries no structural information.
For counterfactual analysis, one needs a method that recovers the reward: AIRL \citep{fu2018} disentangles reward from dynamics, and MCE-IRL \citep{ziebart2010} recovers reward weights via feature matching.
GAIL's role in the econirl pipeline is as a strong imitation learning baseline that exploits MDP structure without requiring feature specification.

\paragraph{Code mapping.}
The econirl implementation at \texttt{src/econirl/contrib/gail.py} provides \texttt{GAILEstimator} with \texttt{GAILConfig}.
Key parameters: \texttt{discriminator\_type} (``tabular'' or ``linear''), \texttt{discriminator\_lr} (default 0.01), \texttt{max\_rounds} (default 100), \texttt{reward\_transform} (``softplus'' or ``logit'').
The tabular discriminator stores $D(s,a)$ as a matrix of logits updated by binary cross-entropy.
The reward signal $\tilde{r} = \text{softplus}(\text{logit}) = \log(1 + e^{\text{logit}})$ is passed to soft value iteration (tabular) or PPO (neural) for the policy step.
Standard errors are available only via bootstrap because there is no structural likelihood.

\bibliography{references}

\end{document}
